%% it/sections/appendici/app_storica.tex
%% ============================================================
%% app_storica.tex — Approfondimento storico + Glossario EN
%% ============================================================

\chapter{Approfondimento Storico}
\markboth{Approfondimento Storico}{Approfondimento Storico}

\section{Cronologia Essenziale}

\begin{tabular}{lp{10cm}}
\toprule
\cinzel Anno & \cinzel Evento \\
\midrule
1194  & Federico II nasce a Jesi. La Sicilia normanda entra nella sfera sveva. \\
1250  & Morte di Federico II. Fine del «Stupor Mundi». \\
1266  & Carlo d'Angiò sconfigge Manfredi a Benevento. \\
1268  & Esecuzione di Corradino a Napoli. Fine della dinastia sveva. \\
1267  & Trattato di Viterbo: Carlo ottiene diritti su Costantinopoli. \\
1277  & Carlo acquista il titolo di re di Gerusalemme. \\
1281  & Elezione di Martino IV, papa filo-angioino. \\
1281  & Carlo pianifica la spedizione contro Bisanzio: 300 navi, primavera 1282. \\
1282 gen & Le requisizioni per la spedizione raggiungono livelli insostenibili. \\
1282 feb & Giovanni da Procida è (forse) a Palermo. La rete si attiva. \\
1282 mar 30 & Vespri Siciliani: rivolta a Palermo, chiesa di Santo Spirito. \\
1282 mar 31 & La rivolta si estende. Quasi tutti i francesi a Palermo sono uccisi. \\
1282 ago & Pietro III d'Aragona sbarca in Sicilia. \\
1282--1302 & Guerra del Vespro. \\
1285  & Morte di Carlo d'Angiò. \\
1302  & Pace di Caltabellotta. La Sicilia resta aragonese. \\
\bottomrule
\end{tabular}

\secsep

\section{Fonti Primarie}

\subsection{Bartolomeo di Neocastro}
\textit{Historia Sicula} (1293 ca.). Scritto da un contemporaneo
siciliano, è la fonte più vicina agli eventi. Presenta una Sicilia
in rivolta giusta contro un tiranno straniero. Va letta consapevoli
del suo punto di vista.

\subsection{Giovanni Villani}
\textit{Nuova Cronica} (XIV sec.). Fiorentino, scrive decenni dopo.
Enfatizza il ruolo di Giovanni da Procida come architetto della
cospirazione. La sua versione è la più drammatica e la più influente
sulla memora collettiva posteriore.

\subsection{Saba Malaspina}
\textit{Rerum Sicularum Historia}. Punto di vista filo-angioino.
Prezioso per capire come Carlo d'Angiò e i suoi sostenitori
interpretarono la rivolta.

\subsection{Il \textit{Liber Iurium} di Genova}
Documenti commerciali e diplomatici che mostrano come la rivolta
venisse percepita nei centri commerciali del Mediterraneo.

\secsep

\section{Storiografia Moderna}

\begin{description}[leftmargin=2em, style=nextline]
  \item[Steven Runciman, \textit{The Sicilian Vespers} (1958)]
    La ricostruzione più completa e leggibile. Accetta in larga
    misura la narrativa della cospirazione aragonese. Ancora
    il punto di partenza obbligato.
  \item[David Abulafia, \textit{Frederick II} (1988)]
    Non dedicato specificamente ai Vespri, ma fondamentale per
    capire la Sicilia sveva e il contesto da cui emerge la rivolta.
  \item[Jean Dunbabin, \textit{Charles I of Anjou} (1998)]
    Il miglior studio su Carlo d'Angiò in inglese. Ridimensiona
    la «cattiveria» di Carlo a favore di una lettura strutturale.
  \item[Henri Bresc, \textit{Un monde méditerranéen} (1986)]
    Studio fondamentale sulla Sicilia medievale come spazio
    economico e culturale.
\end{description}

\secsep

%% ============================================================
%% GLOSSARIO (versione inglese)
%% ============================================================


\filettodoppio